\begin{center}
\SourceLeaf{040}\pdfbookmark[2]{22.}{article-022}22.
\end{center}
Per differentiationem formulae IV prodit (spectando $\psi$ ut quantitatem constantem)
\[
\begin{aligned}
\frac{\mathrm{d}u}{u}&=\frac{1}{2}\left(\operatorname{tang}\frac{1}{2}(v+\psi)-\operatorname{tang}\frac{1}{2}(v-\psi)\right)\mathrm{d}v\\
&=\frac{r\operatorname{tang}\psi}{p}\,\mathrm{d}v;
\end{aligned}
\]
hinc
\[
r r\,\mathrm{d}v=\frac{p r}{u\operatorname{tang}\psi}\,\mathrm{d}u
\]
sive, substituendo pro $r$ valorem ex X,
\[
r r\,\mathrm{d}v=b b\operatorname{tang}\psi\left\{\frac{1}{2}e\left(1+\frac{1}{u u}\right)-\frac{1}{u}\right\}\mathrm{d}u.
\]
Integrando deinde ita, ut integrale in perihelio evanescat, fit
\[
\begin{aligned}
\int r r\,\mathrm{d}v&=b b\operatorname{tang}\psi\left\{\frac{1}{2}e\left(u-\frac{1}{u}\right)-\log u\right\}\\
&=k t\BareRoot{p}\, .\,\BareRoot{(1+\mu)}=k t\operatorname{tang}\psi\, .\,\BareRoot{b}\, .\,\BareRoot{(1+\mu)}.
\end{aligned}
\]
Logarithmus hic est hyperbolicus; quodsi logarithmos e systemate \textsc{Brigg}ico vel generaliter e systemate cuius modulus $=\lambda$ adhibere placet, massaque $\mu$ (quam pro corpore in hyperbola incedente haud determinabilem esse supponere possumus) negligitur, aequatio hancce formam induit:
\[
\mathrm{XI.}\quad\begin{gathered}
\frac{1}{2}\lambda e\frac{u u-1}{u}-\log u=\frac{\lambda k t}{b^{\frac{3}{2}}},\\
\text{sive introducendo }F\\
\lambda e\operatorname{tang}F-\log\operatorname{tang}(45^0+\frac{1}{2}F)=\frac{\lambda k t}{b^{\frac{3}{2}}}\cdot
\end{gathered}
\]
Si logarithmos \textsc{Brigg}icos adhiberi supponimus, habemus \\
$\log\lambda=9{,}637\,7843\,113$, $\log\lambda k=7{,}873\,3657\,527$, sed praecisionem aliquantulum maiorem attingere licet, si logarithmi hyperbolici immediate applicantur. Tangentium logarithmi hyperbolici in pluribus tabularum collectionibus reperiuntur, e.\ g. in iis quas \textsc{Schulze} curavit, maiorique adhuc extensione in \textsc{Beni. Ursin}i Magno Canone Triangulorum Logarithmico, Colon. 1624, ubi per singula $10^{\prime\prime}$ progrediuntur. — Ceterum formula XI ostendit, valoribus reciprocis ipsius $u$, sive valoribus oppositis ipsius $F$ et $v$, respondere valores oppositos ipsius $t$, quapropter partes hyperbolae aequales a perihelioque utrimque aequidistantes temporibus aequalibus describentur.
